\documentclass[10pt]{article}

\usepackage[utf8]{inputenc}

\usepackage[T1]{fontenc}

\usepackage{graphicx}

\usepackage[export]{adjustbox}

\graphicspath{ {./images/} }

\usepackage{amsmath}

\usepackage{amsfonts}

\usepackage{amssymb}

\usepackage[version=4]{mhchem}

\usepackage{stmaryrd}

\usepackage{bbold}

\usepackage{mathrsfs}



\begin{document}

\includegraphics[max width=\textwidth, center]{2023_07_08_ae455b7a522926e1cbfdg-1}

ны м ростком - геометрический объект, возникающий при нахождении квазиклассических приближений с комплексными фазами (см. [1]). Пусть в $2 n$-мерном фазовом пространстве $M^{2 n}$ с симплектич. структурой $\omega^{2}$ задано изотропное $k$-мерное подмногообразие $\Lambda:\left.\omega^{2}\right|_{\Lambda}=0$. Ком лексным ростком на $\Lambda$ называется отображение $r^{n}: q \rightarrow r^{n}(q)$, сопоставляющее для любого $q \in \Lambda$ подпространство $r^{n}(q)$ комплексификации $T_{q}^{\mathbb{C}} M^{2 n}$ касательного к $M$ в точке $q$ пространства, обладаюшее свойствами:



a) лагранжевости $r^{n}(q): \omega^{2}(\xi, \eta)=0$ для любых $\xi, \eta \in r^{n}(q)$;



б) $r^{n}(q) \supset T_{q}^{\mathbb{C}} \Lambda$



в) диссипативности $r^{n}(q): \omega^{2}(\xi, \xi) / i>0$ для любых $\xi \in r^{n}(q) \backslash T_{q}^{\mathbb{C}} \Lambda$.



Векторное расслоение с базой $\Lambda$ и слоями $r^{n}(q)$ называется изотропным подмногообразием с комплексным ростком.



См. также Квантование изотропных торов.



Лит.: [1] Масло в В. П., Комплексный метод ВКБ в нелинейных уравнениях, М., 1977; [2] В ал инь о Б., Доброхотов С. Ю., Нехорошев Н. Н., «Успехи математич. наук», 1984, т. 39, в. 3, c. 233-34.



С. Ю. Доброхотов.



ИЗОTPOПНЫЙ BEKTOP - см. Четырехмерный вектор.



ИЗРАИЛЯ ТЕОРЕМА - см. Черная дыра; теоремы единственности.



ИМПЕДАНС - см. Пассивная система.



ИМПРИМИТИВНАЯ ГPУПกА - группа $G$ взаимно однозначных отображений на себя (подстановок) некоторого множества $S$, для которой существует разбиение множества $S_{\text {в объединение непересекаюшихся подмножеств }} S_{1}, \ldots, S_{m}$, $m \geqslant 2$, обладающее следующими свойствами: число элементов хотя бы в одном из $S_{i}$ больше единицы; для любой подстановки $g \in G$ и любого номера $i, 1 \leqslant i \leqslant m$, существует такой номер $j, 1 \leqslant j \leqslant m$, что $g$ отображает $S_{i}$ на $S_{j}$.



Набор подмножеств $S_{1}, \ldots, S_{m}$ называется с исте мой и м п р и м и т и в ност и, а сами подмножества $S_{i}-$ о бл а стями импримитивности группы $G$. Не импримитивная группа подстановок называется п р и м и т и в н о й.



Понятие И. г. подстановок имеет аналог для групп линейных преобразований векторных пространств. Именно, линейное представление $\rho$ группы $G$ называется и м при м и т и в н ы м, если существует разложение пространства $V$ представления $\rho$ в прямую сумму собственных подпространств $V_{1}, \ldots, V_{m}$, обладающее следующим свойством: для любого элемента $g \in G$ и любого номера $i, 1 \leqslant i \leqslant m$, найдется такой номер $j, 1 \leqslant j \leqslant m$, что $\rho(g)\left(V_{i}\right)=V_{j}$. Если $V$ не обладает разложением указанного типа, то представление $\rho$ называется п р и м и т и в ным. Импримитивное представление $\rho$ называется тран з и т и но и м пр и м и т в ным, если для любой пары подпространств $V_{i}$ и $V_{j}$ из системы импримитивности сушествует такой элемент $g \in G$, что $\rho(g)\left(V_{i}\right)=V_{j}$.



Лит.: [1] Холл М., Теория групп, пер. с англ., М., 1962; [2] Кэртис Ч., Райнер И., Теория представлений конечных групп и ассоциативных алгебр, пер. с англ., М., 1969.



Н. Н. Вильямс, В. Л. Попов.



импульс - то же, что количество движения. См. также Обобщенный импульс.



иМпульС силы - мера действия силы за некоторый промежуток времени; равняется произведению среднего значения силы $\boldsymbol{F}_{\mathrm{cp}}$ на время $t_{1}$ ее действия:



\[

\boldsymbol{S}=\boldsymbol{F}_{\mathrm{cp}} t_{1} \text {. }

\]



И. с. - величина векторная, и направлен он так же, как $\boldsymbol{F}_{\mathrm{cp}}$. Точное значение И.с. за промежуток времени $t_{1}$ определяется интегралом:



\[

S=\int_{0}^{t_{1}} F d t

\]



\textbackslash title\{



212 И30TPONHOE:



При движении материальной точки под действием силы $F$ ее количество движения получает за время $t_{1}$ прирашение, равное И. с.:



\[

S=m v_{1}-m v_{0}

\]



$\left(m v_{0}\right.$ и $m v_{1}-$ соответственно количество движения точки в начале и в конце промежутка времени $t_{1}$ ).



Понятие И. с. широко используется в механике, в частности в теории удара, где величина, равная импульсу ударной силы $F_{\text {уд }}$ за время удара $\tau$, называется ударным имп ульсом.



ИМПУЛЬСНОЕ ПРЕДСТАВЛЕНИЕ К В а Т О во Й М хан и к ( $($ - предста вле н и е) - описание квантовомеханических систем, основанное на разложении векторов состояния $|\psi(t)\rangle$ по базисным векторам $\left|\boldsymbol{p}_{1}, \boldsymbol{p}_{2}, \ldots\right\rangle$, отвечающим определенным значениям импульсов $p_{1}, p_{2}, \ldots$ каждой из частиц. Если число частиц $n$ фиксировано, то



\[

|\psi(t)\rangle=\int d p_{1} d p_{2} \ldots d p_{n}\left|p_{1}, p_{2}, \ldots, p_{n}\right\rangle \times\left\langle p_{1}, p_{2}, \ldots, p_{n} \mid \psi(t)\right\rangle,

\]



где амплитуда $\left\langle p_{1}, p_{2}, \ldots, p_{n} \mid \psi(t)\right\rangle$ представляет собой $n$-частичную волновую функцию в И. П. Вероятность того, что в момент времени $t$ импульс 1-й частицы лежит в интервале $\left(p_{1}, p_{1}+d p_{1}\right)$, импульс 2-й частицы - в интервале $\left(p_{2}, p_{2}+d p_{2}\right)$ и т. д., пропорциональна



\[

\left|\left\langle\boldsymbol{p}_{1}, \boldsymbol{p}_{2}, \ldots, \boldsymbol{p}_{n} \mid \psi(t)\right\rangle\right|^{2} d \boldsymbol{p}_{1} d \boldsymbol{p}_{2} \ldots d \boldsymbol{p}_{n}

\]



Если взаимодействие в системе зависит лишь от относительных расстояний между частицами и отсутствуют внешние поля, нарушающие однородность пространства, то полный импульс $\boldsymbol{P}=\boldsymbol{p}_{1}+\boldsymbol{p}_{2}+\ldots+\boldsymbol{p}_{n}$ сохраняется, и его можно обратить в 0 , переходя в систему центра масс частиц. В результате число независимых импульсов, от к-рых зависит волновая функция, уменьшается на единицу.



Наиболее важное и адекватное применение И. п. находит в квантовомеханич. теории рассеяния. Особенно возрастает роль И. п. при переходе к релятивистскому описанию взаимодействий частиц в квантовой теории поля, где оно объединяется с энергетич. представлением в рамках одного четырехмерного $p$-представления.



См. также Шредингера представление. В. Г. Кадышевский. ИМПУЛЬСНОЕ ПРОСТРАНСТВО - пространство, точкИ которого определяют значения импульсов структурных элементов (частиц) системы. В общем случае - пространство обобщенных импульсов - переменных, канонически сопряженных обобщенным координатам. Размерность И. п. равна полному числу обобщенных координат, то есть числу степеней свободы $S$. Так, для системы $N$ частиц без внутренних степеней свободы размерность И. п. $S=3 N$.



И. п. является подпространством, образующим вместе с пространством обобщенных координат фазовое пространство системы. При классич. описании (замкнутой) системы c $S$ степенями свободы каждое состояние системы в любой момент времени полностью определяется значением $S$ обобщенных координат $q_{i}$ и $S$ обобщенных импульсов $p_{i}$, то есть задается определенной точкой в фазовом пространстве. Соответственно каждая точка И. П. однозначно фиксирует импульсы составляющих систему частиц. В квантовой механике, согласно неопределенностей соотношению, частицы не могут характеризоваться одновременно точно определенными значениями координат и импульсов. Поэтому имеет смысл говорить только о числе состояний в данном (малом) объеме фазового пространства вокруг точки с координатами $\left\{q_{i}, p_{i}\right\}$.



А. Э. Мейерович. ИНВАРИАНТНАЯ МЕРА - для измеримого отображения $T$ измеримого пространства $(\Omega, \mathscr{F})$ в себя, такая мера $\mu$, заданная на $\mathscr{F}$, что



\[

\mu(A)=\mu\left(T^{-1} A\right)

\]



для любого $A \in \mathscr{F}$. Инвариантность меры для $T$ означает также, что



\[

\int f(x) d \mu(x)=\int f(T x) d \mu(x) .

\]





\end{document}